\cotefile{31}
\scanpage{2}
% Faint pencil numeral, upper field, left of title; attribution uncertain.
4
\uline{$\mathrm{cd}^{!}_{p}(k)$ \unsure{pour} $k$ de car. $p$.}

(topolog. fppf)

% archive sheet number: 1

\scanpage{3}
$k$ corps de car $p>0$

On désigne par $\mathrm{cd}_{p}(k)$ la dimension cohomologique, calculée pour les faisceaux de $p$-torsion pour la topologie finie \deleted{\unsure{étale}} \inserted{fid. plat} sur $k$. Regardons d'abord la $\mathrm{cd}'_{p}(k)$, \unsure{relative aux des|relative aux|calculée avec les} groupes algébriques \inserted{(\unsure{comm})} finis sur $k$, \inserted{d'ordre une puiss. de $p$.} Tout tel groupe algébrique se dévisse en
\[
G'\subset G''\subset G
\]
où $G''=G^{0}$ donc $G/G''$ est un groupe étale, et où $G'$ est \inserted{(radiciel)} de t. m. enfin $G''\neq G'$ \inserted{(radiciel)} de type additif. Donc

\uline{Lemme} 1 \quad $\mathrm{cd}'_{p}(k)\leq n\ \Longleftrightarrow\ H^{i}(k,G)=0$ pour $i>n$, lorsque $G$ est \inserted{un $p$-groupe fini sur $k$ de} l'un des types suivants :

a) \deleted{$\unsure{U|F}$}$G$ étale sur $k$, \quad [i.e. \deleted{$\unsure{n|\Longrightarrow}$} $\mathrm{cd}^{\text{ét}}_{p}(k)\leq n$]

b) $G$ \deleted{\unsure{de}}radiciel de t. m. sur $k$

c) $G$ \unsure{est} \deleted{\unsure{ra}}radiciel de t. additif.

% archive sheet number: 2

\scanpage{4}
% Typescript face; read after a half turn. The gaps for signs are visible on the sheet.
est toujours vérifiée si $\deleted{\unsure{T}}G$ est un tore). Alors la démonstration donnée montre que dans l'énoncé de 4.6. \inserted{et les corollaires 4.7. et 4.8. ,} on peut se borner à considérer les sous-groupes de la forme $ {}_{m^{r}}\!G$, avec $r\quad 0$. D'ailleurs, on voit aussitôt \deleted{que} \inserted{lorsque $G$ est \uline{simple} sur $S$, alors} pour tout point $s\quad S$, il existe un voisinage ouvert $U$ de $s$ \deleted{tel} et un entier $m\quad 0$, premier à toutes les caractéristiques \deleted{\unsure{rési|rés} \unsure{résiduelle}} résiduelles de $U$, et qui satisfait les conditions précédentes pour $G\quad U$. (Prendre par exemp\deleted{\unsure{o|b}}le pour $m$ l'ordre de la torsion du type de $\deleted{\unsure{H}}G$ \deleted{dans} en $s$). Sous ces conditions, on trouve des $ {}_{m^{r}}\!G$ qui sont finis et \uline{étales} sur $S$, et dont la famille est schématiquement dense, à condition de se restreindre au dessus de $U$. Cette remarque permet dans certains cas (notamment \deleted{dans l'application} \inserted{ceux faisant intervenir} des théorèmes 3.2. et 3.2. bis, \unsure{cf} \inserted{Appendice -} mais non celle des théo\deleted{\unsure{èr}}rèmes 3.6. et 3.6. bis) de se passer d\deleted{\unsure{e}}u \deleted{\unsure{l}}théorè\deleted{\unsure{e}}me 3.1. faisant intervenir la cohomologie de Hochschild. \deleted{Bien entendu, la remarque analogue s'applique au deux corollaires 4.7. et 4.8.}

5. \uline{Homomorphismes centraux des groupes de type multiplicatif}.

\uline{Théorème} 5.1. Soient $u,v:H\quad G$ deux homomorphismes de $S$-préschémas en groupes, avec $H$ de type multiplicatif et de type fini; \deleted{On suppose} $\deleted{\unsure{H}}G$ de présentation finie sur $S$, \deleted{\unsure{ou|on} $S$ localement noethérien\unsure{.}} Soit $s\quad S$ tel que $u_{s}=v_{s}$, et supposons $u$ \uline{central}. Alors il existe un voisinage ouvert $U$ de $s$ tel que $u_{U}=v_{U}$.

Distinguons les deux cas: a) $S$ localement noethérien, auquel cas on peut se borner à supposer $H$ localement de type fini sur $S$. Cette dernière hypothèse et EGA IV 8 nous permettent de remplacer $S$ par le spectre de $\underline{O}_{S,s}$, donc de supposer $S$ local de point fermé $s$. Quitte à remplacer l'anneau local $A$ par son complété, ce qui introduit un changement de base fidèlement plat et quasi-compact, on peut même supposer $A$ complet. En vertu de 3.4. on aura $u_{n}=v_{n}$ pour tout $n$,

% archive sheet number: 3

\scanpage{5}
Noter que $\mathrm{cd}_{p}^{\text{ét}}(k)$ est égal à $0$ ou $1$, suivant que \deleted{toute} extension \inserted{finie} étale de $k$ est d'ordre premier à $p$, ou non.

Si $G$ est de t. m., alors $G$ est dual d'un groupe fini étale. Ce dernier est quotient d'un module galoisien, \unsure{libre} de rang fini sur $\mathbb{Z}$, donc $G$ est contenu dans un tore, \unsure{donc on a} une suite exacte
\[
0\longrightarrow G\longrightarrow T\longrightarrow T'\longrightarrow 0
\]
$T$ et $T'$ des tores. D'où le

\uline{Lemme} 2 \quad Pour que $\mathrm{cd}_{p}^{\text{t.m.}}(k)$ soit $\leq n$, il f. et s. que les conditions suivantes soient vérifiées : pour tt tore $T$ sur $k$, $H^{i}(k,T)(p)$ est nul pour $i>n$, et $p$-divisible si $i=n$. [Ou encore : pour tte extension finie \unsure{séparable} $k'$ de $k$ [il suffit d'ordre premier à $p$] \deleted{$\mathrm{Br}(k)(p)$} $H^{i}(k',G_{m})=0$
% margin note written sideways; paired loop carries on the square bracket.
\margin{si $i>n$, et $H^{n}(k',G_{m})$ est $p$-divisible]}

Comme on a $0\longmapsto {}_{p}T\longrightarrow T\xrightarrow{p}T\longrightarrow 0$

d'où
\[
\deleted{\left\{ {}_{p}H^{\unsure{i|\prime}}(T)=\operatorname{Im} H^{i}({}_{p}T)\right.}
\]

% archive sheet number: 4

\scanpage{6}
ou comme d'habitude l'indice $n$ indique la réduction mod \deleted{la puissance} $\underline{m}^{n+1}$ ($\underline{m}$ étant l'idéal maximal de $A$). Pour tout entier $m\quad 0$, dénotant par $ {}_mu,{}_mv$ les homomorphismes $ {}_mH\quad G$ induits par $u,v$, on aura donc aussi $({}_mu)_n=({}_mv)_n$. Ceci étant vrai pour tout $n$, et $ {}_mH$ étant \uline{fini} sur $S$, en vertu de 2.2. , il s'ensuit qu'on aura $ {}_mu={}_mv$. Ceci étant vrai pour tout $m$, on aura $u=v$ en vertu de 4.7. a).

b) Cas général. Nous \deleted{supposerons ici} pouvons supposer $S$ \deleted{\unsure{local de point fermé $s$}} \inserted{affine, $S=\operatorname{Spec}(A)$.} Quitte à faire une extension fidèlement plate et quasi compacte de la base, on peut supposer de plus $H$ diagonalisable, donc de la forme $D_S(M)$, avec $M$ un groupe \deleted{\unsure{de}}commutatif de type fini. Introduisons comme dans la démonstration de 4.6. la famille filtrante croissante des sous-$\underline{Z}$-algèbres \deleted{$A_i$} de type fini $A_i$ de $A$. Comme $\deleted{H}\unsure{G|H}$ est \unsure{ce|de} présentation finie sur $S$, il provient d'un \deleted{schéma} préschéma en groupes de présentation finie $\deleted{H}G_i$ sur un $S_i$ . D'ailleurs, $H$ provient de $H_i=D_{S_i}(M)$ De plus, pour $\deleted{i}j$ grand, en vertu de EGA IV 8, les homomorphismes \deleted{$u,v$} $u,v$ proviennent d'homomorphismes de $S_j$-groupes $\deleted{\unsure{H}}u_j,v_j:H_j\quad G_j$, \deleted{Soit $s_j$ im} et $u_j$ sera central pour $j$ grand. D'ailleurs, si $s_j$ est l'image de $s$ dans $S_j$, on a $u_{j\,s_j}=v_{j\,s_j}$ . On peut alors appliquer a) à la situation sur $S_j$, d'où aussitôt la conclusion annoncée .

\uline{Cor\deleted{\unsure{a}}ollaire} 5.2. Soient $u:H\quad G$ un homomorphisme de $S$\deleted{(}-préschémas en groupes, avec $H$ de type multiplicatif et de type fini, et $H$ de présentation finie (ou seulement localement de \deleted{pr} type fini sur $S$, lorsque $S$ est localement noethérien). Soit $s\quad S$ et supposons que $u_s$ soit l'homomorphisme unité. Alors il existe un voisinage ouvert $U$ de $s$ tel que $u_U$ soit l'homomorphisme unité.

On prouve de la même façon les résultats suivants (qui contiennent d'ailleurs les précédents, comme on voit en introduisant \deleted{\unsure{l}}\unsure{d|l}es \deleted{\unsure{grais gra|grism gra}} graphes des homomorphismes en groupes précédents):
% archive sheet number: 5

\scanpage{7}
\[
\begin{aligned}
H^i({}_pT)&\longrightarrow {}_pH^i(T)\longrightarrow 0\\
0&\longrightarrow H^i(T)_p\longrightarrow H^{i+1}({}_pT)
\end{aligned}
\]
on voit que les conditions sont nécessaires.

Pour la suffisance, on tire d'abord le fait que $H^{\deleted{\unsure{n}}i}(G)_{\deleted{\unsure{p}} \unsure{p}}=$ pour $i\geq n+2$, reste à vérifier $H^{n+1}(G)=0$ donc que $H^n(T)\deleted{{}_{\unsure{(p)}}}\longrightarrow H^n(T')\deleted{{}_{\unsure{(p)}}}$ est surjectif. \deleted{\unsure{avec|Bien} \unsure{T} \unsure{De fait. Faisons} Le conoyau est un \unsure{groupe} de $p$-torsion, donc on peut se borner à regarder $H^n(T)_{(p)}\longrightarrow H^n(T')_{(p)}$. C'est une \unsure{bonne|horreur} que $G$ soit $p$-divisible,}

Noter que si $k'/k$ est une extension finie séparable de $k$, alors pour tt $G'$ sur $k'$, on a
\[
H^i(k,\prod_{k'/k}G')\simeq H^i(k,G),
\]
ce qui prouve que l'hypothèse faite sur les tores sur $k$ est \uline{stable} par extension de base séparable. Soit $k'$ une telle extension de degré premier à $p$, alors l'hom trace $\prod_{k'/k}G_{k'}\longrightarrow G$ est \deleted{surjectif} \inserted{un épim.}
% archive sheet number: 6

\scanpage{8}
\deleted{\uline{Théorème} 5.1. bis \quad Soient $G$ un $S$-préschéma en groupes de présentation finie, $H_1$ et $H_2$ deux sous-préschémas en groupes de type multiplicatif, \deleted{\unsure{sur $S$ tel que $H$|s $S$ tel que $H$}} Soit $s\quad S$ \deleted{\unsure{et su}} \deleted{\unsure{y}}tel que $H_{1\,s}=H_{2\,s}$, et supposons $H_1$ central. Alors il existe un voisinage ouvert $U$ de $s$ tel que $H_{1\,U}=H_{2\,U}$.}

\deleted{\uline{Corol}laire 5.2. bis}

\uline{Corollaire} 5.3. \quad Soient $u,H,G$, comme dans 5.2. Supposons de plus $G$ \uline{séparé} sur $S$. Soit $U$ l'ensemble des $s\quad S$ tels que $u_s$ soit l'homomorphisme unité. Alors $U$ est une partie à la fois ouverte et fermée de $S$, et $u_U:H_U\quad G_U$ est l'homomorphisme unité.

Il reste à prouver seulement que $U$ est fermé. Or \deleted{pour tout entier} $n\quad 0$, soit $ {}_nu:{}_nH\quad G$ l'homomorphisme induit, et $U_n$ l'ensemble des $s\quad S$ tels que $ {}_nu_s$ soit l'homomorphisme unité. En vertu de 4.6. \unsure{appliqué|applique} aux fibres $H_s$, $U$ est l'intersection des $U_n$, donc on est ramené à prou\deleted{ver que ces derniers sont fermés, donc on peut supposer $G$ \uline{fini} sur $S$.} soit $K=\operatorname{Ker}u$, comme $H$ est séparé sur $S$, c'est un sous-préschéma \uline{fermé} de $H$, et $U$ est l'ensemble des $s\quad S$ tels que $K_s=$\deleted{\unsure{encore}}$H_s$ . On voit alors \deleted{\unsure{facistem|facislem|facilemen}} faci\deleted{\unsure{m}}lement que $U$ est fermé, par exemple comme application de Exp V,6.1. ,ou en \deleted{\unsure{se ramenant}} notant qu'on peut supposer $S$ réduit et $U$ dense dans $S$, donc schématiquement dense dans $S$, ce qui implique $H_U$ schématiquement dense dans $H$ puisque $H$ est plat sur $S$, et comme $u$ et l'homomorphisme unité $e$ coïncident sur $H_U$, ils coincident sur $H$.

La démonstration est la même que pour 5.1. , en rempla,cant 3.4. par 3.4. bis et 4.7. a) par 4.7. b).

\deleted{\uline{Corollaire} 5.3. bis}
% archive sheet number: 7

\scanpage{9}
[ \uline{NB}. $\prod_{k'/k}G_{k'}$ est de t. m. ] donc grâce à $H^{n+2}=0$ sur les groupes de t. m. on voit que
\diagram{\begin{tikzcd}[ampersand replacement=\&,row sep=0.25em]
H^{n+1}(\prod_{k'/k}G_{k'})\arrow[r] \& H^{n+1}(G)\quad\text{est surjectif,}\\
\rotatebox{90}{$\simeq$}\\
H^{n+1}(G_{k'})
\end{tikzcd}}
donc pour vérifier que $H^{n+1}(G)=0$, il suffit de le faire pour $H^{n+1}(G_{k'})$. Cela nous ramène au cas où $G$ est ``\unsure{splitté|splité}'' par une $p$-extension, \uline{ce qui implique que} $G$ \uline{se dévisse en groupes isomorphes à} $\mu_p$. Donc on est ramené à vérifier que $H^{n+1}(\mu_p)=0$, ce qui se voit sur la suite exacte
\[
0\longrightarrow\mu_p\longrightarrow G_m\longrightarrow G_m\longrightarrow 0
\]

\uline{Lemme} 3 \deleted{\unsure{S}} On a
\[
\operatorname{cd}^{\substack{\rm \unsure{plat|rad}\\\rm add}}_p(k)=
\begin{cases}0&\text{si }k\text{ parfait}\\1&\text{si }k\text{ imparfait}\end{cases}
\]

En effet, un groupe radiciel additif se dévisse en groupes $\alpha_p$ \unsure{\ldots\,.}
% archive sheet number: 8

\scanpage{10}
Généralités sur les \unsure{proobjets|pro objets|produits} adiques....

\unsure{4 p}
% archive sheet number: 9

\scanpage{11}
\uline{Exemple} \qquad $A$ anneau (

$f\in A$ élément $A$-régulier

$g\in A$ élément $(A/fA\inserted{{}=B})$-régulier, \deleted{\unsure{S}}

on suppose $A/fA+gA\neq 0$ i.e $B/gB\neq 0$

[On peut supposer $A$ noethérien régulier de dim 2, $A/fA$ noeth. régulier de dim 1, $A/fA+gA$ noeth régulier de dim 2, \unsure{p. ex|p. ex.} $A$ anneau local noeth de dim 2, $(f,g)$ système de paramètres ]

On va construire un module $M$ sur $A$ t.q. que
\[
M=\bigcup_{n\geq 0}{}_{f^n}M,\qquad {}_{f^{n+1}}M/{}_{f^{\unsure{n+0|n}}}M\simeq A/fA \tag*{(i)}
\]
[en particulier, $M$ est de support $\subset V(f)$ et est $fA$-cofini]

\[
M/fM\simeq(A/fA+gA)^{(\mathbb N)}=(B/gB)^{(\unsure{\mathcal N|N|\mathbb{N}})}\quad\text{donc}\tag*{(ii)}
\]
$M/fM$ n'est pas $fA$-cofini.

\uline{Construction} On construit par récurrence les $ {}_{f^n}M$ comme $E_n$, l'hypothèse de récurrence étant
\[
\left\{\begin{aligned}
\mathrm{a)}_n\;&f^nE_n=0\\
\mathrm{b)}_n\;&{}_{f^{k+1}}E_n/{}_{f^k}E_n\simeq A/fA=B\qquad(0\leq k\leq n-1)\\
\mathrm{c)}\;&E_n/fE_n\simeq\deleted{A/fA} B+(B/gB)^{n-1}
\end{aligned}\right.
\]
\margin{\textcircled{\unsure{M|4}}.}
% archive sheet number: 10

\scanpage{12}
\mspage{\handwritten{30}}
% A single long oblique stroke crosses the opening through the theorem heading.
\deleted{\deleted{on a nécessairement} puisque ses composantes connexes sont ouvertes ($S$ étant localement noethérien). Mais alors on a nécessairement $Z=J_S$, où $J$ est une partie de $I$, et on a \deleted{\unsure{$Z=$|$Z=K$}} donc $U=\emptyset$ ou $U=S$, suivant que $J$ est infini ou fini, ce qui donne la conclusion voulue.

\uline{Corollaire} 5.1\deleted{\unsure{3}}\handwritten{4}. Soi\deleted{t}\inserted{ent} $S$ un préschéma localement noethérien et \uline{normal}, $\deleted{\unsure{R}}\handwritten{P}$ un $S$-préschéma \deleted{localement} constant t\deleted{\unsure{r}}o\deleted{\unsure{o}}rdu quasi-isotrivial. Alors les composantes connexes de $\deleted{R}\handwritten{P}$ sont finies sur $S$.

Comme $\deleted{R}\handwritten{P}$ est étale sur $S$, on sait que $\deleted{R}\handwritten{P}$ est normal (SGA I 9.5.) donc ses composantes \deleted{irréductibles} connexes sont ses composantes irréductibles. On peut évidemment supposer $S$ connexe donc irréductible, soit $s$ son point générique. Comme $R$ est plat sur $S$, ses points maximaux, i.e. les points\handwritten{\unsure{'}} génériques $x_i$ des composantes irréductibles $R_i$ de $R$, sont les points maximaux de $R_s$, et la fibre \deleted{de $R$} gén\handwritten{é}rique $R_{i\,s}$ de $R_i$ est la composante irréductible de $R_s$\handwritten{,} adhérence de $x_i$ dans $R_s$. Or $R_s$ étant étale sur $k(s)$, est somme disjointe de spectres d\deleted{\unsure{e}}'extensions finies séparables de $k(s)$, donc \deleted{les} pour tout $i$, $R_{i\,s}$ est fini (en fait, réduit à un seul point). Donc 5.1\deleted{2}\handwritten{3}. s'applique et montre que les $R_i$ sont finis sur $S$.

\uline{Remarque} 5.1\deleted{\unsure{4}}\handwritten{5}. Le raisonnement s'applique, plus généralement, si on suppose seulement $S$ géométriquement unibranche au lieu de normal, i.e. que le morphisme \deleted{\unsure{S}}canonique $\overset{\handwritten{\sim}}{S}\handwritten{\longrightarrow}S$, où $\overset{\handwritten{\sim}}{S}$ est le normalisé de $S_{\text{réd}}$\handwritten{,}\deleted{\unsure{,}} est radiciel. Donc le théorème 5.1\deleted{5}\handwritten{6}. ci-dessous est valable également sous cette condition plus générale.

\uline{Théorème} 5.1\deleted{5}\handwritten{6}. Soit $S$ un préschéma localement noethérien normal. Alors tout groupe de type multiplicatif} \deleted{\unsure{et de type fini|de type fini} \unsure{présentation|quasi-isotrivial}} \inserted{et de type fini} sur $S$\deleted{\unsure{, est et en parti|, et et en partic}} est \uline{isotrivial}.

On peut supposer en effet \inserted{$S$ connexe, donc} $H$ de type constant $M$. Il suffit d'appliquer 5.13 . avec $P=\underline{\operatorname{Isom}}_{S\text{-gr}}(H,G)$, où $G=D_S(M)$\handwritten{,}
\handwritten{en raisonnant comme dans 5.1\unsure{2|1}. (ii) $\Longrightarrow$ (i); \deleted{\illegible{2}} \inserted{On peut aussi} appliquer \unsure{à|ou} \deleted{\unsure{avec}} 5.11. \deleted{\illegible{4}}, en utilisant 5.14. pour $R=D_S(\underline{H})$.}
\margin{\handwritten{5.11. \textcircled{\unsure{M|M$^0$}}.}}
% archive sheet number: 11

\scanpage{13}
la deuxième \deleted{\unsure{à|se}} se ramène à
\[
E_n={}_{f^n}E_{n+1}
\]
qui par $\mathrm a_n)$ équivaut à $\supset$, donc à ceci

Soient $x\in E_n$, $h\in A$, \deleted{$f^nx$.} tels qu'il existe $\lambda\in A$, avec
\[
f^n(x,h)=\lambda(\xi_n,f)
\]
alors on a $h\in fA$. En effet, on aura $f^nh=\lambda f$ d'où $\lambda=f^{n-1}h$, (\uline{car $f$ est $A$-régulier}) puis $f^nx=\lambda\xi_n$ \deleted{\unsure{i.e.} $f^{\unsure{n}}$.} i.e.
\[
0=f^{n-1}h\xi_n,\quad\text{or}\quad (\overbrace{f^{n-1}\xi_n}^{y_n})\in {}_fE_n\quad\text{par }\mathrm a_n),
\]
\deleted{\unsure{d'où}} et $ {}_fE_n\simeq A/fA=B$ par $\mathrm b_n)$, \deleted{\unsure{donc}} enfin par construction on a $\xi_n\notin{}_{f^{n-1}}E_n$ i.e. \deleted{$f^{n}$} $y_n\neq0$, donc $hy_n=0\Longrightarrow h\in fA$, cqfd.

Enfin, vérifions $\mathrm c_n$, on a
\[
\begin{aligned}
E_{n+1}/fE_{n+1}&\simeq(E_n\times A)/f(E_n\times A)+A(\xi_n,f)\\
&\simeq(E_n/fE_n\times A/fA)/A(\dot{\xi}_n,\dot f)\\
&\simeq\bigl([B\times(B/gB)^{n-1}]\times B\bigr)/B(g,0,0)\\
&\simeq B/gB\times(B/gB)^{n-1}\times B\simeq B\times(B/gB)^n.
\end{aligned}
\]
% archive sheet number: 12

\scanpage{14}
\mspage{\handwritten{29}}
% One continuous oblique stroke crosses the typescript, from the head to the penultimate line.
\deleted{images dans $S$ sont des parties $S_i$ à la fois ouvertes et fermées, et elles recouvrent $S$. Si $S$ est connexe ou quasi-compact, il existe donc un ensemble fini d'indices \deleted{tel} $i$ tels que les $S_i$ recouvrent $S$, soit $S'$ la réunion des $P_i$ correspondants, \deleted{qui est donc tel que} \inserted{Alors} $S'\handwritten{\mathord{-}}S$ \deleted{\unsure{r}}est fini étale surjectif, et posant $P'=P\times_S S'=\underline{\operatorname{Isom}}_{S'\text{-gr}}(H',G')$, on voit que $P'$ a une section sur $S'$, donc il existe un isomorphisme de $S'$-groupes $H'=H\times_S S'\handwritten{\xrightarrow{\sim}}G'=G\times_S S'=D_{S'}(M)$, ce qui prouve que $S'$ splitte $H$. Le cas où \deleted{\unsure{S}}les composantes \deleted{irréductibles}\inserted{connexes} de $S$ sont \deleted{\unsure{connex}} ouvertes se ramène aussitôt au cas où $S$ est connexe. \quad - \quad Remarquons d'ailleurs qu'on peut, lorsque $H$ est de type fini sur $S$,\deleted{\unsure{connexe|connex}} et de type constant, donner \deleted{un}le critère d'isotrivialité suivant (où il n'est plus nécessaire de faire aucune restriction sur $S$): $R$ est réunion d'une \uline{suite} de parties à la fois ouvertes et fermées finies sur $S$.

\uline{Lemme} 5.1\deleted{2}\handwritten{3}. Soient $S$ un préschéma localement noethérien et connexe, $\deleted{R}\handwritten{P}$ un $S$-préschéma constant tordu quasi-isotrivial, $Z$ une partie à la fois ouverte et fermée de $\deleted{R}\handwritten{P}$, telle qu'il existe un $s\handwritten{\in}S$ tel que $\deleted{\unsure{R}} \handwritten{Z}_s$ soit fini. \margin{\handwritten{$Z_s$}} Alors $\deleted{\unsure{R}}\handwritten{Z}$ est fini sur $S$. \margin{\handwritten{$Z$}}

\deleted{Oublions un instant que $S$ est supposé connexe, et} \inserted{Ne supposons pas d'abord $S$ connexe;} soit $U$ l'ensemble des $s\handwritten{\in}S$ tels que $Z_s$ soit fini, nous allons prouver que $U$ est à la fois ouvert et fermé, et que $Z\handwritten{|}U$ est fini sur $U$. C'est là évidemment un \deleted{énoncé} énoncé essentiellement équivalent à 5.1\deleted{2}\handwritten{3}. mais qui a l'avantage d'être ``de nature locale'' sur $S$ pour la topologie étale (disons), ce qui nous ramène au cas où $\deleted{R}\handwritten{P}$ est constant, i.e. de la forme $\deleted{\unsure{K}}I_S$, où $I$ est un ensemble. (NB l'hypothèse localement noethérienne n'est pas perdue par un changement de base étale;} c'est là où on utilise la quasi-isotrivialité de $\deleted{R}\handwritten{P}$ sur $S$). On peut de plus supposer $S$ connexe\handwritten{,} \deleted{Mais alors}
% archive sheet number: 13

\scanpage{15}
On construit $E_{n+1}$ comme extension de $A/fA$ par $E_n$, \deleted{\illegible{1}} déduite de l'extension \unsure{standart|standard} de $A$ de $A/fA$ par $A$ par un hom $\deleted{\unsure{E}}A\longrightarrow E_n$ défini par un élément $\xi_n$ de $E_n$ (l'extension $E_{n+1}$ ne dépend que de $\xi_n$ mod $fE_n$).
\diagram{\begin{tikzcd}[ampersand replacement=\&]
0\arrow[r,dash] \& A\arrow[r,"f"]\arrow[d,"\xi"] \& A\arrow[r]\arrow[d] \& A/fA\arrow[r,dash]\arrow[d,Rightarrow] \& 0\\
0\arrow[r,dash] \& E_n\arrow[r] \& E_{n+1}\arrow[r] \& A/fA\arrow[r,dash] \& 0
\end{tikzcd}}

Notons que \inserted{par $\mathrm a_n$} $fE_n\subset\deleted{f}{}_{f^{n-1}}E_n$ donc
\diagram{\begin{tikzcd}[ampersand replacement=\&,column sep=1.3em,row sep=0.05em]
0\arrow[r,dash] \& {}_{f^{n-1}}E_n/fE_n\arrow[r] \& E_n/fE_n\arrow[r] \& E_n/{}_{f^{n-1}}E_n\arrow[r,dash] \& 0\\
\&\&\&\rotatebox{90}{$\simeq$}\\
\&\&\& A/fA=B
\end{tikzcd}}
\unsure{Comme|donc} $g\longrightarrow$ \unsure{injectif} dans $B$, \deleted{et} \deleted{\unsure{les seuls}} \deleted{\unsure{f}} \deleted{\unsure{hom de}} \deleted{\unsure{les}} \deleted{\unsure{formes}} \inserted{hom} $B$-linéaires de $B+(B/gB)^{n-1}$ dans $B$ sont \unsure{de de|ts de} 2 la \deleted{première} forme $\lambda\,\mathrm{pr}_1$ ($\lambda\in B$), et s'il \deleted{\unsure{est}} est surjectif\unsure{?}, $\lambda$ est inv. Donc l'isom $\mathrm c)$ \deleted{\unsure{est}} \inserted{peut être pris} compatible avec la structure d'extension précédente [où $ {}_{f^{n-1}}E_n/fE_n\simeq(B/gB)^{n-1}$)

Nous choisirons $\dot\xi_n=(\deleted{\unsure{\lambda,\ }}\dot g,0)$. \deleted{Alors}

Vérifions les propriétés $\mathrm a_{n+1})$, $\mathrm b_{n+1})$, $\mathrm c_{n+1})$.

La première résulte \unsure{trivialement} de $(\mathrm a_n)$ et $fA/fA=0$
% archive sheet number: 14

\scanpage{16}
\mspage{\handwritten{28}}
% Oblique line from head of sheet to the bottom correction.
\deleted{\uline{Lemme} 5.1\deleted{\unsure{1}}\handwritten{2}. Soient $S$ un préschéma, $R$ un préschéma constant t\deleted{\unsure{r}}o\deleted{\unsure{o}}rdu sur $S$. Alors tout sous-préschéma fermé \deleted{\illegible{1}} \inserted{\handwritten{$Z$}} de $R$ qui est \deleted{de type fini}\inserted{quasi-compact} sur $S$, est fini sur $S$.

En effet, on est ramené au cas où $R$ est constant, donc de la forme $I_S$, où $I$ est un ensemble, donc réunion filtrante croissante des $J_S$, où $J$ parcourt les parties finies de $I$. On peut supposer de plus $S$ affine, alors \deleted{\unsure{$R_i$}} $Z$ est quasi-compact, donc contenu dans un des $J_S$. Comme $J_S$ est fini sur $S$, il en est de même de $Z$.

Le lemme 5.1\deleted{\unsure{1}}\handwritten{2}. établit déjà l'assertion entre parenthèses dans 5.\deleted{10}\handwritten{11}. (ii). \deleted{L'équivalence (ii) \quad (iii) est} \inserted{L'implication (ii)$\handwritten{\Longrightarrow}$(iii), et l\unsure{1|l}implication inverse} dans le cas $S$ localement noethérien est alors claire. Prouvons que (i)$\handwritten{\Longrightarrow}$(ii), pour ceci, soit $S'\handwritten{\longrightarrow}S$ un morphisme étale fini surjectif qui splitte $H$ donc $R$, de sorte que $R'$ s'écrit \deleted{comme} $M_S=\deleted{\unsure{R_m}}\handwritten{\coprod}_{m\handwritten{\in}M}R'_m$, où \deleted{\unsure{sont|tous|tout}} les $R'_m$ sont des ouverts disjoints de $R'$, \inserted{$S'$-}isomorphes à $S'$. Soit $g:R'\handwritten{\to}R$ la projection, qui est un morphisme fini étale surjectif, \inserted{donc un morphisme ouvert et fermé\unsure{;};} alors les $R_m=g(R'_m)$ sont des parties à \deleted{a}la fois ouvertes et fermées de $R$, et évidemment quasi-compacts sur $S$ puisque les $R_m$ le sont. \deleted{Enfin, prouvons (ii) \quad (i)}

Enfin, prouvons (ii)$\handwritten{\Longrightarrow}$(i) lorsque $H$ est de type fini sur $S$\handwritten{,} \inserted{\handwritten{$S$ quasi compact ou connexe.}} Soit $M$ son type, on peut écrire $M=\deleted{\unsure{T_{\oplus}|Z_{\oplus}}}\underline Z^r\times N$, où $r$ est un entier $\handwritten{\geq}0$ et $N$ un groupe fini. Considérons le préschéma
\[
P=\underline{\operatorname{Isom}}_{S\text{-gr}}(H,G)\handwritten{\subset}\underline{\operatorname{Hom}}_{S\text{-gr}}(H,G)\handwritten{\unsure{=|\Rightarrow}}Q
\]
où $G=D_S(M)$, cf 5.\handwritten{8}. et 5.\handwritten{10}. \deleted{Alors $Q$ est isomorphe} On a un isomorphis
\[
Q\handwritten{\simeq}\underline{\operatorname{Hom}}_{S\text{-gr}}(G,\underline G_m^r)\times\underline{\operatorname{Hom}}_{S\text{-gr}}(G,N)\handwritten{\simeq}R^r\times E\quad,
\]
où $E=\underline{\operatorname{Hom}}_{S\text{-gr}}(G,N)$ est fini sur $S$. Il s'ensuit que $Q$ est \deleted{somme de}\inserted{réunion de sous-}préschémas à la fois ouverts et fermés $Q_i$ finis sur $S$. \deleted{De même, Il s'ensuit que} \inserted{\handwritten{Donc}} $P$ est réunion \deleted{\unsure{d'un|disj}} de \deleted{\unsure{présc|prés}} sous-préschémas à la fois ouverts et fermés $P_i$\inserted{\handwritten{$=P\cap Q_i$}}} finis sur $S$. \deleted{L}Comme ils sont étales sur $S$, leurs
% archive sheet number: 15

\scanpage{17}
De plus, la construction donne un hom can
\[
E_n\xrightarrow{u_n}E_{n+1}
\]
tel que l'hom correspondant
\diagram{\begin{tikzcd}[ampersand replacement=\&,row sep=0.3em]
u_n\otimes_A B: \& E_n/fE_n\arrow[r,dash]\arrow[d,dash] \& E_{n+1}/fE_{n+1}\arrow[d,dash]\\
\& Q_n \& Q_{n+1}
\end{tikzcd}}
soit \uline{nul} sur $g(E_n/fE_n)$, et induise un hom injectif de \inserted{\unsure{module libre|un mod libre}} $\deleted{E_n/fE_n}Q_n\otimes_B B/gB$ dans \deleted{\unsure{et définissant une base libre|est définie et un hom injec} $Q_{n+1}$ \unsure{en appliquant une base de}} et de $Q_{n+1}\otimes B/gB$, appliquant une base sur une base ce sont des modules libres).

Donc
\[
\begin{aligned}
\varinjlim_n Q_n&=\varinjlim_n Q_n\otimes_B B/gB\simeq\varinjlim (B/gB)^{n-1}\\
&\simeq (B/gB)^{(\mathbb N)}.
\end{aligned}
\]
Prenons donc
\[
M=\varinjlim_n E_n
\]
Alors $M$ satisfait aux conditions voulues.
% archive sheet number: 16

